% 中文摘要（通识课程论文不要求英文摘要，已删除）
\cleardoublepage
\begin{center}
    \bfseries \zihao{3} 摘~要
\end{center}

\vskip 6pt

几十年来MRI的空间分辨率一直在提高，但更能说明这门技术变化的，是被编码的对象不断更换：从组织形态，到水分子的扩散，到血流灌注和血管通透性，再到脑的功能活动。本文以信号产生与空间编码为共同基础，沿三条线索展开：扩散加权成像及其向扩散张量、扩散峰度成像的延伸；动态对比增强磁共振从强化形态观察走向药代动力学定量；血氧水平依赖功能磁共振从任务态定位走向静息态脑网络分析。结合急性脑梗死、前列腺癌、乳腺癌、抑郁症和卒中康复的案例，讨论这些技术在疾病检出、病理评估、疗效观察与风险预测中各自走到了哪一步。我的基本判断是：MRI能够支持临床决策，靠的并不是序列和算法数量的堆积，而是让某一类信号与一个具体的临床问题之间建立起可检验的对应关系，并用规范的采集、增量价值的评价和可重复性检验来约束这种对应。文章最后从生物医学工程专业学生的角度，谈技术灵敏度、参数可解释性与临床可实施性之间如何取舍。

\vskip 12mm

\noindent\textbf{关键词：}磁共振成像；扩散加权成像；动态对比增强；功能磁共振；多参数成像；临床决策
